StaticKD keeps \val{statickd.test.pct_teacher} of the teacher macro-F1 on the unified test set and \val{statickd.dev.pct_teacher} on the development set, at an inference cost two to three orders of magnitude lower (\val{eff.statickd.speedup_onnx}$\times$ against ONNX int8 and \val{eff.statickd.speedup_pt}$\times$ against PyTorch). This agrees with the finding that order-free representations are already competitive for topic classification~\cite{galke2022bow}. Names of institutions, sports terms, and medical vocabulary carry most of the topic signal. The remaining gap concentrates in two places. The first is labels whose vocabulary overlaps with other topics. The largest gaps to the teacher occur for \val{label.biggest_gap}. The hardest labels for StaticKD, given here with their F1, are \val{label.hardest}. They are also among the hardest for the teacher. \emph{Society} and \emph{human interest} are known to be difficult for human annotators as well~\cite{kuzman2025llm}. The second is languages that are underrepresented in the distillation data. StaticKD reaches \val{grp.distil_tier.cukup.pct} of the teacher score in languages with sufficient distillation data but only \val{grp.distil_tier.tidakada.pct} in languages without any.

With the same distillation data, MiniLM-L6 closes part of this gap (StaticKD $-$ MiniLM \val{sig.test.minilm_kd.delta}, $p$ \val{sig.test.minilm_kd.p}), whereas DistilmBERT does not (\val{sig.test.distilmbert_kd.delta}, $p$ \val{sig.test.distilmbert_kd.p}). Both cost one to two orders of magnitude more than StaticKD. MiniLM comes from a distillation of XLM-R-large and DistilmBERT from mBERT, which suggests that the quality of the initial multilingual representation matters more than the number of layers.

The ablations show where the knowledge of the student comes from (RQ3). The largest source is the multilingual distillation data labeled by the teacher. Without it the test macro-F1 drops by \val{abl.emm_only.test_macro.delta}, consistent with Sparse Distillation~\cite{ye2022sparse}, which also relies on a large transfer set. The cross-lingually aligned initialization comes next, and the compact variant gives the most direct evidence for it. In Indian languages that hardly occur in the distillation data (for example bn, pa, or), most tokens are never updated during training and keep their initial potion embeddings. StaticKD still classifies these languages with a mean accuracy of \val{lang.vlowcov.static_acc}, yet the accuracy falls to \val{lang.vlowcov.compact_acc} once the tokenizer pruning removes those tokens. Transfer to languages without data is therefore carried by embedding rows that were never trained but live in the same space as the trained rows, so the linear head learned on data-rich languages applies to them too.

The frozen-embedding ablation shows the other side of the same mechanism. Fine-tuning moves the embeddings of tokens from languages in the distillation data and improves in-distribution fit. Freezing them lowers the development score by \val{abl.frozen.dev_macro.delta}. Tokens of other languages stay in the original potion space, where the head that adapted to the moved space fits them less well. In languages with no or little distillation data, frozen embeddings are better by \val{ablgrp.frozen.tidakada.delta} and \val{ablgrp.frozen.sedikit.delta}. This is representation drift caused by fine-tuning part of the vocabulary. The final recipe keeps fine-tuning because the recipe was selected on a development set in the training languages. For deployments that target languages without distillation data, frozen embeddings or a regularizer that keeps the embeddings close to potion are the safer choice, and both still fold into a table.

Soft labels, by contrast, add little or nothing. Hard teacher labels do not differ significantly from soft labels at $\tau=1$ (development \val{abl.hard.dev_macro.delta}, test \val{abl.hard.test_macro.delta}). Raising the temperature to $\tau=2$ lowers the held-out agreement without a significant change on the test set, and $\tau=4$ lowers every metric. This pattern fits the capacity-gap literature~\cite{zhang2025law,wang2026rckd}. A plausible explanation is that one linear layer over averaged embeddings cannot represent the inter-class structure in the tail of the teacher distribution (dark knowledge)~\cite{zhao2022dkd}. A sharp distribution ($\tau=1$) or hard labels give an attainable target, whereas a high $\tau$ asks the student to imitate a distribution shape beyond its capacity. For very-low-capacity students, effort is better spent on the coverage of the transfer set (languages, sites, topics) and on aligned initialization than on shaping the soft labels.

Per-language performance is driven mainly by how well a language is represented in the distillation and pre-training data, not by its script (Kruskal--Wallis $p$ \val{fac.kw.p}). The weakest languages are African languages outside XLM-R pre-training that are also missing from the distillation data (\val{lang.outxlmr.langs}). The teacher still reaches an accuracy of \val{lang.outxlmr.teacher_acc} on them, perhaps thanks to subwords shared with other languages and to sentence context, while StaticKD reaches only \val{lang.outxlmr.static_acc}. The cheapest fix needs no change of architecture. News in these languages can simply be added to the transfer set and labeled by the teacher. The data ablation and the correlation between agreement and the amount of data ($\rho=\val{fac.data.rho}$) support this step.

Some variants move agreement with the teacher and accuracy on editorial labels in opposite directions. Frozen embeddings lower the held-out agreement (\val{abl.frozen.heldout_acc.delta}) but raise the test macro-F1 (\val{abl.frozen.test_macro.delta}). Cutting the data to 130 thousand documents lowers the agreement (\val{abl.data130k.heldout_acc.delta}) without a significant change on the test set (\val{abl.data130k.test_macro.delta}). In line with~\cite{stanton2021kd}, students should be selected with labels that do not come from the teacher and with test data outside the training distribution, not with fidelity alone.

Test scores should be read according to the nature of the labels. Two thirds of the labels are editorial topics, the place where the newsroom put an article, translated to IPTC. An article about a criminal case involving an athlete can rightly sit in a sports section while a model labels it \emph{crime}. Both answers are plausible under IPTC, which allows several topics per article. Absolute scores on editorial labels thus measure agreement with the categorization of publishers. Comparisons between models remain fair because all models are scored against the same labels, and the model ranking does not change on the subset with capped site counts.

Compared with Sparse Distillation~\cite{ye2022sparse}, StaticKD replaces an n-gram table trained from scratch with pre-trained multilingual embeddings folded into 17 logits per token. The model becomes far smaller and transfers to languages that are missing from the transfer set. Inference is the recurring cost in a news pipeline, since every incoming article is classified, and one million articles per day need \val{eff.teacher_onnx_int8.cores1M} cores with the ONNX int8 teacher against \val{eff.statickd.cores1M} with StaticKD (Section~\ref{sec:res-eff}). Int8 quantization of the teacher only gives a constant-factor speed-up.

For a real-time pipeline on servers without a GPU, StaticKD processes one million articles per day with \val{eff.statickd.cores1M} cores, a \val{final.size}~MB file, and \texttt{numpy} and \texttt{tokenizers} as its only dependencies. Its probabilities are better calibrated than those of the teacher (test ECE \val{ece.statickd.test} against \val{ece.teacher.test}), which makes a cascade possible. With $\theta=0.8$, the cascade accuracy (\val{casc.08.accuracy}) equals that of the teacher alone (\val{casc.101.accuracy}) at \val{casc.08.docs_per_sec} articles per second per core, against \val{eff.teacher_onnx_int8.docs_per_sec_1c} for the teacher alone. This mode suits batch jobs that need teacher-level accuracy. The compact variant lowers RAM but is only advisable when all target languages are covered by the distillation data.
